\documentclass[11pt]{article}
\usepackage[utf8]{inputenc}
\usepackage[T1]{fontenc}
% es-noshorthands como el libro: con los atajos de babel activos, "->" de TikZ rompe.
\usepackage[spanish,es-noshorthands]{babel}
\usepackage{tikz}
\usetikzlibrary{positioning,calc}
\usepackage{etoolbox}

% Externalización: sólo si el driver (main-ext.tex) define \LvExternalize antes de \documentclass.
% Se externaliza únicamente lo que está dentro de un entorno figure: el tikzpicture con
% "remember picture, overlay" de más abajo, externalizado, desaparecería del PDF (plan.md,
% tabla de la sección 2, fila 11).
\ifdefined\LvExternalize
  \usetikzlibrary{external}
  \tikzexternalize[prefix=figuras-cache/]
  \tikzexternaldisable
  \AtBeginEnvironment{figure}{\tikzexternalenable}
  \AtEndEnvironment{figure}{\tikzexternaldisable}
\fi

\begin{document}

\section{Figuras}

Primera figura, dentro de \texttt{figure}, que es la que se externaliza.

\begin{figure}[ht]
  \centering
  \begin{tikzpicture}
    \draw[->] (0,0) -- (3,0) node[right] {$\omega$};
    \draw[->] (0,0) -- (0,2) node[above] {$|H|$};
    \draw[thick,blue,domain=0:3,samples=40] plot (\x,{2/sqrt(1+\x*\x*\x*\x)});
  \end{tikzpicture}
  \caption{Respuesta de un filtro de segundo orden.}
\end{figure}

Segunda figura, con otro dibujo, para que haya más de una en el archivo y el orden importe en el nombre del cache.

\begin{figure}[ht]
  \centering
  \begin{tikzpicture}
    \foreach \k in {0,...,5} {
      \draw[fill=gray!30] ({cos(60*\k)},{sin(60*\k)}) circle (0.1);
    }
    \draw (0,0) circle (1);
  \end{tikzpicture}
  \caption{Seis polos sobre una circunferencia.}
\end{figure}

\section{Fuera de figure}

Un \texttt{tikzpicture} con \texttt{remember picture, overlay} \emph{fuera} de \texttt{figure}, como los íconos y portadillas del template del libro. Si algo lo externaliza, desaparece del PDF.

\begin{tikzpicture}[remember picture,overlay]
  \node[anchor=north east,red] at (current page.north east) {marca de agua};
\end{tikzpicture}

\end{document}
